\chapter*{Lời Mở Đầu}
\addcontentsline{toc}{chapter}{Lời Mở Đầu}
\markboth{Lời Mở Đầu}{Lời Mở Đầu}

\begin{truyen}{Lời Mở Đầu}{Opening Words}
\chuhoa{Q}{uyển này đi từ một điều rất quen trong lớp học: học sinh hỏi rất nhiều về chuyện học, nhưng không phải câu nào cũng được hỏi thành tiếng.}
\textit{(This volume begins from something very familiar in classrooms: students have many questions about learning, but not every question is spoken aloud.)}

Có câu các em hỏi thẳng: “Vì sao phải học?”, “Điểm thấp có phải thất bại không?”, “Em học chậm có phải dốt không?”
\textit{(Some questions are direct: “Why do we have to study?”, “Does a low score mean failure?”, “If I learn slowly, am I stupid?”)}

Cũng có những câu không được nói ra nhưng hiện trong ánh mắt, trong cách ngại giơ tay, trong nỗi sợ kiểm tra, trong sự nản khi so với bạn bè, và trong những lúc muốn bỏ cuộc.
\textit{(Other questions are never spoken, but appear in the eyes, in the fear of raising a hand, in anxiety before tests, in discouragement after comparison, and in moments of wanting to give up.)}

Quyển này gom 100 câu hỏi rất gần với đầu óc học sinh. Không nhằm biến việc học thành một bài giảng đạo đức dài, mà để trả lời ngắn, rõ, thật, và đủ gần với những điều các em đang sống trong trường lớp.
\textit{(This book gathers 100 questions that stay close to the student mind. It does not aim to turn learning into a long moral lecture, but to answer briefly, clearly, honestly, and close to school life as students actually live it.)}

\end{truyen}

\begin{baihoc}
Học không chỉ là chuyện điểm số. Học còn là chuyện hiểu mình, hiểu đời, bớt sợ sai, và lớn lên đàng hoàng hơn.
\textit{(Learning is not only about scores. It is also about understanding oneself, understanding life, fearing mistakes less, and growing into a more grounded person.)}
\end{baihoc}

\begin{ghinhoanh}
\textbf{Short English to remember:}

Learning is not only about scores.

It is also about growing into a steadier person.

\end{ghinhoanh}

\begin{tuhocnhanh}
\textbf{Gợi ý để dạy và tự nghĩ / Teaching and reflection prompts}

1. Học sinh đang hỏi gì về việc học mà mình chưa nghe ra? \textit{What questions about learning are students asking that we have not yet heard clearly?}

2. Khi nói về học, mình đang nói quá nhiều về kết quả hay đủ về con người? \textit{When we speak about learning, are we talking too much about results and not enough about the person?}

3. Mình muốn học sinh đi hết quyển này với điều gì còn lại trong lòng? \textit{What do I want students to carry in their hearts after finishing this volume?}
\end{tuhocnhanh}

\ngancach
